% Generated draft: requires mathematical and source review.
\documentclass[11pt]{article}
\usepackage[margin=25mm]{geometry}
\usepackage{fontspec}
\setmainfont{Latin Modern Roman}
\usepackage{amsmath,amssymb,mathtools,unicode-math}
\setmathfont{Latin Modern Math}
\usepackage{enumitem,needspace,xurl,hyperref}
\setlength{\parindent}{0pt}
\setlength{\parskip}{5pt}
\setlength{\emergencystretch}{3em}
\newcommand{\N}{\mathbb{N}}
\newcommand{\Z}{\mathbb{Z}}
\newcommand{\Q}{\mathbb{Q}}
\newcommand{\R}{\mathbb{R}}
\newcommand{\C}{\mathbb{C}}
\newcommand{\F}{\mathbb{F}}
\newcommand{\A}{\mathbb{A}}
\newcommand{\PP}{\mathbb{P}}
\newcommand{\E}{\mathbb{E}}
\newcommand{\eps}{\varepsilon}
\newcommand{\dd}{\,\mathrm{d}}
\begin{document}
\section*{UnsolvedMath ID: 30001029}
\textbf{Source problem code:} catalogue entry
\subsection{Definitions and mathematical statement}
In conventional notation, a Hamiltonian isotopy on a symplectic manifold $(M,\omega)$ is a smooth path $\phi_t$ generated by time-dependent Hamiltonians $H_t$:
\[
\frac{d}{dt}\phi_t=X_{H_t}\circ\phi_t,
\qquad \iota_{X_{H_t}}\omega=-dH_t,
\qquad t\in[0,1].
\]
The sign convention is immaterial to the meaning of being Hamiltonian; support, completeness and boundary conditions must be fixed in the intended setting.

The source conjecture is that Hamiltonian isotopies are naturally strongly intersecting in the precise sense used for its symplectic homoclinic Floer invariant. Thus, for every Hamiltonian isotopy admissible in that construction, the required naturally-strongly-intersecting condition is conjectured to hold. The phrase is a technical condition, not a substitute for an ordinary assertion that two sets intersect.

For context, a homoclinic point of a hyperbolic fixed point $p$ of a diffeomorphism $f$ lies in $W^s_f(p)\cap W^u_f(p)$ and is distinct from $p$, where stable and unstable sets are defined by convergence to $p$ under forward and backward iterates respectively. This describes the terminology only. The source does not specify the manifold class, starting map, fixed-point data, branches of stable and unstable manifolds, their continuation under isotopy, or the actual definition of natural strong intersection. In particular, one must not simply require every map of an identity-based Hamiltonian path itself to have a hyperbolic fixed point; the construction may instead deform a specified dynamical map. That setup is not supplied.

Consequently a self-contained universal formula for the intersection predicate cannot be recovered here. No transversality condition, intersection number, pairwise branch requirement or Floer-invariance theorem is invented in its place. The requested conjecture is retained with its source-dependent qualification.

The corrected-verified note says the report explicitly states this conjecture. The supplied literature assessment checked 2026-08-21 labels it open and says a later primary paper still calls Hamiltonian isotopies naturally strongly intersecting only conjecturally, with no proof located. These are dated attributed claims.
\subsection{Short English statement}
Retain the conjecture that Hamiltonian isotopies satisfy the natural strong-intersection condition needed for the report's homoclinic Floer construction. The precise dynamical setup and technical intersection predicate are missing and must be recovered for a self-contained statement.
\subsection{Sources}
\begin{itemize}
\item Local UnsolvedMath registry, permanent ID 30001029, source code catalogue entry; supplied complete record and corrected-verified report-check note. Literature checked 2026-08-21: the assessment is supplied dated curation, not independently verified here.
\item Geometric Group Theory, Hyperbolic Dynamics and Symplectic Geometry (2009), Oberwolfach Reports, DOI 10.4171/owr/2008/42, pp. 2337--2406: \url{https://doi.org/10.4171/owr/2008/42}. Proposed-year metadata gives 2008.
\item Homoclinic points and Floer homology: \url{https://arxiv.org/abs/1704.02403}.
\end{itemize}
\end{document}
